\begin{explanationbox}
	\textbf{\textcolor{CExplanation}{一句话直觉}}
	集合的包含关系就是逻辑推出的方向。范围小的条件可以推出范围大的条件，反之不必然。

	\medskip
	\textbf{\textcolor{CExplanation}{核心拆解}}
	\begin{description}[style=nextline, leftmargin=2.8em, labelsep=0.5em, itemsep=0.45em]
		\item[\textbf{要点一（集合与逻辑对应）：}]
			每个条件对应一个解集，逻辑关系“若$p$则$q$”等价于解集包含关系$A\subseteq B$。
		\item[\textbf{要点二（包含方向定充分）：}]
			若$A\subseteq B$，则$x\in A\Rightarrow x\in B$，所以$p\Rightarrow q$，即$p$是$q$的充分条件。
		\item[\textbf{要点三（特殊情况分类）：}]
			$A=B$时互推（充要），$A\subsetneqq B$时充分不必要，$B\subseteq A$时必要。
	\end{description}

	\medskip
	\textbf{\textcolor{CExplanation}{几何本质（包含关系可视化）}}
	在数轴上，集合的包含关系直观表现为区间的覆盖。若$A$的区间被$B$完全覆盖，则$A$中元素必在$B$中，对应逻辑推出。

	\medskip
	\textbf{\textcolor{CExplanation}{代数意义（集合运算转化）}}
	判断$p\Rightarrow q$只需验证$A\subseteq B$，即$A\cap B = A$或$A\cup B = B$，转化为集合的运算结果。

	\medskip
	\textbf{\textcolor{CExplanation}{考点价值（判断充要条件）}}
	\begin{itemize}[leftmargin=*, itemsep=0.5em, topsep=0.4em]
		\item \textbf{考法一（包含关系判充分）：} 给出两个条件，直接比较解集包含关系得出结论。
		\item \textbf{考法二（参数范围求充要）：} 已知充分性或必要性，反向确定参数满足的集合包含关系。
	\end{itemize}

	\medskip
	\textbf{\textcolor{CExplanation}{顿悟点}}
	\textbf{方向是关键：谁包含谁就对应谁推出谁。小集合推出大集合，大集合不能推出小集合。}

	\medskip
	\textbf{\textcolor{CExplanation}{使用场景}}
	\begin{itemize}[leftmargin=*, itemsep=0.5em, topsep=0.4em]
		\item \textbf{场景一（条件转化）：} 将条件转化为解集后直接比较包含关系。
		\item \textbf{场景二（参数问题）：} 由充分或必要条件建立集合包含的不等式。
	\end{itemize}

\end{explanationbox}
